\section*{TD NSI :  Opération "Golden Sac"}

\textbf{Objectif de la mission :} \textit{Infiltration réussie. Vous êtes dans la salle des coffres. Votre sac à dos de transport est limité à \textbf{10 kg}. Vous devez choisir les objets les plus chers pour maximiser votre butin avant que la police n'encercle le bâtiment. Bonne chance, agent.}

\begin{Exercise}[title={Infiltration du système de sécurité }]
La base de données des objets est stockée dans un dictionnaire nommé \texttt{musee}. 
\begin{lstlisting}[style=pythonstyle, language=Python]
# "Nom de l'objet": [poids_en_kg, prix_en_keuros]
musee = {
    "Statue d'or": [15, 200],
    "Tableau de maitre": [5, 150],
    "Vase Ming": [2, 50],
    "Armure de samourai": [25, 100],
    "Diademe": [1, 40]
}
\end{lstlisting}

\begin{enumerate}
    \item \textbf{Vrai ou Faux ?} L'instruction \texttt{musee[0]} permet de récupérer la Statue d'or. \dotfill
    \item \textbf{QCM :} Quelle commande renvoie exactement le prix du \texttt{"Diademe"} (\texttt{40}) ?
    \begin{itemize}
        \item[$\square$] \texttt{musee["Diademe"]} \quad $\square$ \texttt{musee["Diademe"][1]} \quad $\square$ \texttt{musee[1]["Diademe"]}
    \end{itemize}
    \item  Ajoutez une \texttt{"Piece rare"} de \texttt{1}kg valant \texttt{10}k\texteuro{} au dictionnaire : \dotfill
    \item  Un complice tape \texttt{print(musee["vase ming"])}. Pourquoi Python affiche-t-il une erreur \texttt{KeyError} ? \dotfill
    \item Quelle instruction permet de modifier le prix du \texttt{"Vase Ming"} pour le passer à \texttt{60}k\texteuro{} ? \dotfill
\end{enumerate}

\begin{Comprendre}
 \textbf{Indice du Hacker :} 
Dans un dictionnaire de listes, \texttt{dico[cle]} renvoie la liste entière. Pour avoir un élément précis de cette liste, on ajoute un deuxième crochet : \texttt{dico[cle][0]} pour le poids, \texttt{dico[cle][1]} pour le prix.
\end{Comprendre}
\end{Exercise}

\begin{Exercise}[title={Analyse des contraintes }]
Votre sac est limité à \textbf{10 kg}.
\begin{enumerate}
    \item  Quel est le poids total de tous les objets du musée ? Est-ce transportable ? \dotfill
    \item  Si vous prenez le \texttt{Tableau} et le \texttt{Vase}, quel est le poids total utilisé ? \dotfill
    \item Quelle est la valeur (en k\texteuro) de la combinaison : \texttt{Tableau + Vase + Diademe} ? \dotfill
    \item \textbf{Logique :} Pourquoi est-il impossible de prendre la \texttt{Statue d'or} ? \dotfill
    \item  Si vous deviez choisir \textbf{un seul} objet pour avoir le plus d'argent, lequel serait-ce ? \dotfill
\end{enumerate}
\end{Exercise}

\newpage

\begin{Exercise}[title={Simulation des Stratégies Gloutonnes }]
Complétez les deux tableaux pour comparer les résultats des stratégies.

 \textbf{Stratégie A : "Priorité à la Valeur" (Le plus cher d'abord)} 
\begin{center}
\begin{tabular}{|c|c|c|c|c|}
\hline
\textbf{Étape} & \textbf{Objet choisi} & \textbf{Prix (k\texteuro)} & \textbf{Poids utilisé} & \textbf{Place restante} \\ \hline
1 & Statue d'or & NON (Trop lourd) & 0 kg & 10 kg \\ \hline
2 & Tableau de maitre & 150 & 5 kg & 5 kg \\ \hline
3 &  &  &  &  \\ \hline
4 &  &  &  &  \\ \hline
5 &  &  &  &  \\ \hline
\end{tabular}
\end{center}

 \textbf{Stratégie B : "Priorité à la Légèreté" (Le plus léger d'abord)}
\begin{center}
\begin{tabular}{|c|c|c|c|c|}
\hline
\textbf{Étape} & \textbf{Objet choisi} & \textbf{Prix (k\texteuro)} & \textbf{Poids utilisé} & \textbf{Place restante} \\ \hline
1 & Diademe & 40 & 1 kg & 9 kg \\ \hline
2 &  &  &  &  \\ \hline
3 &  &  &  &  \\ \hline
4 &  &  &  &  \\ \hline
5 &  &  &  &  \\ \hline
\end{tabular}
\end{center}

\begin{enumerate}
    \item Quelle solution est la meilleure ? \dotfill
\end{enumerate}

\begin{Comprendre}
 \textbf{Rappel :} Un algorithme glouton ne revient jamais en arrière. Si un objet est trop lourd, on passe définitivement au suivant dans la liste des priorités.
\end{Comprendre}
\end{Exercise}


\begin{Exercise}[title={Inspection du code de sélection }]
\begin{lstlisting}[style=pythonstyle, language=Python]
def choisir_meilleur(inventaire, deja_pris, place_restante):
    elu = None
    max_prix = -1
    for nom, infos in inventaire.items():
        poids_obj = infos[0]
        prix_obj = infos[1]
        if nom not in deja_pris and poids_obj <= place_restante:
            if prix_obj > max_prix:
                max_prix = prix_obj
                elu = nom
    return elu
\end{lstlisting}
\begin{enumerate}
    \item Pourquoi initialise-t-on \texttt{max\_prix} à \texttt{-1} plutôt qu'à \texttt{0} ? \dotfill
    \item  À quoi sert la liste \texttt{deja\_pris} dans la condition \texttt{if} ? \dotfill
    \item Que contient la variable \texttt{infos} à chaque tour de la boucle \texttt{for} ? \dotfill
    \item Que renvoie cette fonction si aucun objet ne rentre dans le sac ? \dotfill
    \item Si on voulait la Stratégie B (le plus léger), quelle ligne du code faudrait-il changer ? \dotfill
\end{enumerate}
\end{Exercise}

\begin{Exercise}[title={Le Script de Fuite (Final) }]
\textit{Complétez l'algorithme pour automatiser le remplissage du sac.}

\begin{lstlisting}[style=pythonstyle, language=Python]
def braquage_glouton(dico_musee, limite_sac):
    butin = []            # Objets voles
    poids_actuel = 0      
    fini = False          
    
    while not fini:
        # 1. Calculer l'espace encore disponible
        place_libre = .......... - poids_actuel
        
        # 2. Chercher l'objet optimal (Appel de la fonction de l'exercice 4)
        cible = choisir_meilleur(dico_musee, butin, ..........)
        
        if cible == .......... : 
            # 3. Si aucun objet n'est trouve, on arrete la boucle
            fini = ..........
        else:
            # 4. On ajoute l'objet au butin
            butin.append(..........)
            # 5. On met a jour le poids total du sac
            poids_actuel = poids_actuel + dico_musee[..........][0]
            
    return butin, poids_actuel

# Lancement de l'operation
objets_voles, total_kg = braquage_glouton(musee, 10)
print(f"Extraction terminee ! Objets : {objets_voles} | Poids : {total_kg}kg")
\end{lstlisting}
\end{Exercise}


\begin{Comprendre}
 \textbf{Attention :} \texttt{dico\_musee[cible]} accède à la liste \texttt{[poids, prix]}. N'oubliez pas que pour le poids, l'indice est \texttt{0}.
\end{Comprendre}
